% Build: make focusing   (LuaLaTeX -> PDF -> vector SVG and transparent PNG)
% Exact page: 224.2 x 70 mm. No caption.
\documentclass[tikz,border=0pt]{standalone}
\usepackage{fontspec}
\setsansfont{Arial}
\setmainfont{Arial}
\usepackage{graphicx,amsmath,pgfplots}
\pgfplotsset{compat=1.18}
\usetikzlibrary{arrows.meta,calc}

% Matplotlib tab10: C0/C1 are already used by the original XAS/XANES SVGs.
% All badges use 50% white; internal objects use the saturated base colour.
\definecolor{mplBlue}{HTML}{1F77B4}
\definecolor{mplOrange}{HTML}{FF7F0E}
\definecolor{mplGreen}{HTML}{2CA02C}
\definecolor{mplPurple}{HTML}{9467BD}
\definecolor{mplOlive}{HTML}{BCBD22}
\colorlet{crystalInk}{mplGreen!75!black}
\colorlet{cztInk}{mplPurple!80!black}
\colorlet{simInk}{mplOlive!65!black}

\begin{document}
\begin{tikzpicture}[x=1cm,y=1cm,line cap=round,line join=round,
  every node/.style={font=\sffamily,inner sep=0pt},
  title/.style={font=\sffamily\bfseries\fontsize{12}{14}\selectfont},
  detail/.style={font=\sffamily\fontsize{9}{11}\selectfont,align=center},
  badge/.style={draw=black,line width=.32mm},
  flow/.style={draw=black!85,line width=.42mm,
    -{Stealth[length=2.3mm,width=1.7mm]},shorten <=11.6mm,shorten >=11.6mm},
  axisArrow/.style={black,line width=.20mm,-{Stealth[length=1.0mm,width=.8mm]}},
  mathlabel/.style={font=\fontsize{8}{9}\selectfont}]
\path[use as bounding box] (0,0) rectangle (22.42,7);
\clip (0,0) rectangle (22.42,7);
\begin{pgfinterruptboundingbox}

% Clockwise workflow: model the instrument, select energy, excite the sample,
% detect fluorescence, reconstruct the near-edge spectrum. No feedback claim.
\coordinate (SIM) at (11.21,4.95);
\coordinate (CRYSTAL) at (18.37,4.20);
\coordinate (XAS) at (16.25,1.35);
\coordinate (CZT) at (6.17,1.35);
\coordinate (XANES) at (4.05,4.20);
\draw[flow] (SIM) -- (CRYSTAL);
\draw[flow] (CRYSTAL) -- (XAS);
\draw[flow] (XAS) -- (CZT);
\draw[flow] (CZT) -- (XANES);

% Preserve the user's original vector artwork, with its unused margins cropped.
\node at (XAS) {\includegraphics[width=21.6mm,trim=18mm 18mm 18mm 18mm,clip]{.build/XAS.pdf}};
\node at (XANES) {\includegraphics[width=21.6mm,trim=18mm 18mm 18mm 18mm,clip]{.build/XANES.pdf}};

% CRYSTAL: intrinsic sigma reflectivity of a flat Ge(220) crystal at 10 keV,
% calculated with crystalpy (see README). The bent wafer below is schematic.
% The curve illustrates energy selection; it is not a bent-Ge(880) prediction.
\begin{scope}[shift={(CRYSTAL)}]
  \filldraw[badge,fill=mplGreen!50!white] (0,0) circle (1.06);
  \begin{axis}[at={(-.73cm,.16cm)},anchor=south west,
    width=1.47cm,height=.58cm,scale only axis,
    xmin=-100,xmax=100,ymin=0,ymax=1.05,hide axis]
    \addplot[draw=crystalInk,line width=.31mm]
      table[x index=0,y index=1] {4_Crystal/reflectivity_Ge220.dat};
  \end{axis}
  \draw[axisArrow] (-.76,.14) -- (.83,.14);
  \draw[axisArrow] (-.76,.14) -- (-.76,.75);
  \node[mathlabel] at (-.56,.76) {$\mathcal{R}$};
  \node[mathlabel] at (.64,-.02) {$\Delta\theta$};
  % Broad, dark curved wafer and a few legible lattice planes.
  \path[fill=crystalInk,draw=black,line width=.25mm]
    (-.80,-.26) .. controls (-.30,-.58) and (.30,-.58) .. (.80,-.26)
    -- (.80,-.47) .. controls (.30,-.79) and (-.30,-.79) .. (-.80,-.47) -- cycle;
  \path[fill=mplGreen,draw=black,line width=.25mm]
    (-.80,-.08) .. controls (-.30,-.40) and (.30,-.40) .. (.80,-.08)
    -- (.80,-.26) .. controls (.30,-.58) and (-.30,-.58) .. (-.80,-.26) -- cycle;
  \draw[black,line width=.17mm]
    (-.80,-.17) .. controls (-.30,-.49) and (.30,-.49) .. (.80,-.17);
  \foreach \xx/\drop in {-.64/.09,-.42/.175,-.21/.225,0/.24,.21/.225,.42/.175,.64/.09}{
    \draw[black,line width=.17mm] (\xx,-.08-\drop) -- (\xx,-.26-\drop);
  }
  \node[font=\sffamily\bfseries\fontsize{8}{9}\selectfont] at (0,-.83) {Ge};
\end{scope}

% CZT: a 20-anode array drawn directly in TikZ, then two measured traces.
% The purple sensor and electrodes are schematic; sample values are unchanged.
\begin{scope}[shift={(CZT)}]
  \filldraw[badge,fill=mplPurple!50!white] (0,0) circle (1.06);
  \path[fill=cztInk,draw=black,line width=.20mm]
    (-.79,.51) -- (.71,.51) -- (.80,.61) -- (.80,.39)
    -- (.71,.29) -- (-.79,.29) -- cycle;
  \path[fill=mplPurple,draw=black,line width=.20mm]
    (-.79,.51) -- (.71,.51) -- (.80,.61) -- (-.70,.61) -- cycle;
  \draw[black,line width=.16mm] (.71,.51) -- (.71,.29);
  \foreach \i in {0,...,19}{
    \pgfmathsetmacro{\xx}{-.747+.073*\i}
    \filldraw[fill=mplPurple!25!white,draw=cztInk,line width=.09mm]
      (\xx,.35) rectangle ++(.043,.105);
  }
  \begin{axis}[at={(-.75cm,-.66cm)},anchor=south west,
    width=1.39cm,height=.85cm,scale only axis,
    xmin=0,xmax=1,ymin=-1.15,ymax=1.15,hide axis]
    \addplot[draw=cztInk,line width=.29mm]
      table[x index=0,y index=1] {5_Pulse_Shape/waveforms.dat};
    \addplot[draw=cztInk,line width=.29mm]
      table[x index=0,y expr=\thisrowno{2}-1.12] {5_Pulse_Shape/waveforms.dat};
  \end{axis}
  \draw[black,line width=.31mm] (.70,.04) -- (.78,-.04) -- (.92,.16);
  \draw[black,line width=.31mm] (.72,-.38) -- (.88,-.54) (.72,-.54) -- (.88,-.38);
\end{scope}

% SIMULATION: a clean ray bundle and detector-response distribution.
% Both are schematic, with no simulated acceptance or resolution implied.
\begin{scope}[shift={(SIM)}]
  \filldraw[badge,fill=mplOlive!50!white] (0,0) circle (1.06);
  % Ray paths are behind the source, optic and detector.
  \foreach \xx/\yy in {-.10/.55,0/.53,.10/.55}{
    \draw[simInk,line width=.24mm]
      (-.65,.12) -- (\xx,\yy) -- (.67,-.03);
  }
  \draw[simInk,line width=.25mm,-{Stealth[length=1.15mm,width=.9mm]}]
    (-.53,.20) -- (-.26,.37);
  \draw[simInk,line width=.25mm,-{Stealth[length=1.15mm,width=.9mm]}]
    (.23,.34) -- (.47,.14);
  \filldraw[fill=simInk,draw=black,line width=.23mm,rounded corners=.35mm]
    (-.85,-.02) rectangle (-.66,.29);
  \fill[fill=mplOlive] (-.69,.075) rectangle (-.64,.19);
  \path[fill=mplOlive,draw=black,line width=.25mm]
    (-.26,.69) .. controls (-.10,.60) and (.10,.60) .. (.26,.69)
    -- (.26,.56) .. controls (.10,.47) and (-.10,.47) .. (-.26,.56) -- cycle;
  \filldraw[fill=simInk,draw=black,line width=.23mm]
    (.65,-.18) rectangle (.84,.14);
  \foreach \yy in {-.12,-.04,.04}{
    \draw[mplOlive!30!white,line width=.15mm] (.68,\yy) -- (.81,\yy);
  }
  \draw[axisArrow] (-.68,-.71) -- (.77,-.71);
  \draw[axisArrow] (-.68,-.71) -- (-.68,-.20);
  \path[fill=mplOlive!85!white]
    (-.61,-.69) -- (-.39,-.69)
    .. controls (-.20,-.69) and (-.22,-.30) .. (-.04,-.29)
    .. controls (.15,-.28) and (.13,-.69) .. (.34,-.69)
    -- (.65,-.69) -- cycle;
  \draw[simInk,line width=.31mm]
    (-.61,-.69) -- (-.39,-.69)
    .. controls (-.20,-.69) and (-.22,-.30) .. (-.04,-.29)
    .. controls (.15,-.28) and (.13,-.69) .. (.34,-.69) -- (.65,-.69);
\end{scope}

% Titles and explanatory labels remain outside the circles.
\node[title] at (11.21,6.70) {SIMULATION};
\node[detail] at (11.21,6.28) {Ray tracing\enspace$\cdot$\enspace Detector response};
\node[title] at (20.75,4.50) {CRYSTAL};
\node[detail] at (20.75,3.89) {Bent Ge optics\\Energy selection};
\node[title] at (18.90,1.65) {XAS};
\node[detail] at (18.90,1.04) {X-ray absorption\\spectroscopy};
\node[title] at (3.20,1.65) {CZT};
\node[detail] at (3.20,1.04) {Multichannel detection\\Pulse-shape analysis};
\node[title] at (1.48,4.50) {XANES};
\node[detail] at (1.48,3.89) {Near-edge\\structure};

% Clear centre: the project name only.
\node[font=\sffamily\bfseries\fontsize{29}{31}\selectfont]
  at (11.21,3.04) {PRISM};
\end{pgfinterruptboundingbox}
\end{tikzpicture}
\end{document}
